% abstract.tex

The hydrogen atom in the A=1 Discrete Causal Lattice framework
(Paper~I~\cite{menendez2026geometry}) is a joint Arnold-tongue
attractor between the proton and electron sessions, lock-in
occurring when their orbital frequency ratio matches the lattice's
$\omega \cdot R_1 = \pi/3$ resonance to within the tongue width
$\Delta\omega_\text{tongue}$.  An external mass $M$ at distance $d$
imposes a differential clock-density gradient across the atom of
order $\Delta\omega_\text{tidal} \sim M \cdot R_1 / d^3$.  When the
tidal detuning exceeds the tongue width, the joint resonance breaks
and the orbit dissolves -- gradient ionization, a third channel
beyond thermal and pressure ionization.  This paper derives the
critical condition $M > \Delta\omega_\text{tongue} \cdot d^3 / R_1$
analytically -- the \emph{quantum Roche limit}, sharing the $d^3$
scaling of the classical Roche limit with $\Delta\omega_\text{tongue}$
playing the role of the classical body density -- and confirms it
numerically via the inherited \texttt{exp\_18\_tidal\_ionization.py}
g-scan, which measures the orbital escape time $T_\text{escape}(g)$
as the imposed gradient $g = M / d^2$ is swept.  The framework's
parameter-free $M_\text{min}(d)$ ionization curve gives a
temperature-independent inner edge for neutral hydrogen near compact
objects, testable against 21\,cm HI absorption morphology and
accretion-disk Fe K$\alpha$ line profiles.  The paper closes with the
atomic ISCO -- the innermost stable atomic orbit, an analogue of the
geodesic ISCO at $r = 6GM/c^2$, below which no neutral matter can
exist regardless of temperature.
